\scanpage{544}
\mspage{455}

\unsure{obtenu|obtenues} ainsi est de cardinal $\leq\operatorname{card} M\leq\pi_0$, donc l'ens des relations d'équivalence $R_i$ $(i\in I)$ est lui aussi de cardinal $\leq\pi_0$. Cela montre, $\unsure{\pi|\pi_0}$ étant \unsure{grd} devant $\pi_0$, qu'il existe une relation d'équivalence \inserted{de la forme} $R_u$ $(u\in M)$, et un ens. cofinal $I'$ de $I$, tel que $\forall i\in I'$, on ait $R_i=R_u$. Cela montre que, quitte
à remplacer $I$ par l'ens ordonné cofinal $I'$, qu'on peut supposer que l'on a
\[
R_i=R_u\qquad\forall i\in I.
\tag*{\makebox[0pt][r]{\makebox[\displaywidth][l]{(4.15.8.5)}}}
\]
Soit alors
\[
\deleted{\unsure{N|V}}\unsure{E|N}=M/R_u=M/R_i\qquad(\forall i\in I),
\]
on voit donc que les $\alpha_{ji}\in M$ \inserted{(pour $i\leq j$)} sont tels que l'application
\[
\tau_{\alpha_{ji}}\colon v\longmapsto\alpha_{ji}v:M\mathbin{\unsure{-|\to}}M
\]
est compatible avec la relation d'équivalence $R_u$, donc passe au quotient en une
application $\overline{\alpha_{ji}}\colon\deleted{N}E\longrightarrow\deleted{N}E$, \deleted{\unsure{et}} laquelle est d'ailleurs \uline{injective}.

Soit $M'$ le sous-monoïde de $M$ formé des $u\in M$ qui ont la propriété \inserted{de stabilité} précédente, \deleted{et $\overline{\deleted{\unsure{N|E}}M'}$ le sous-monoïde} \inserted{ainsi que la propriété d'injectivité} qu'on vient de dire, et $\overline{M'}$ le monoïde \deleted{\unsure{quoti}} image de $M'$ dans le monoïde
\[
\operatorname{Mon}(\deleted{N}E)\subset\operatorname{End}_{\unsure{\mathrm{Ens}|\mathrm{En}}}(\deleted{N}E)
\]
% archive sheet number: 543
